% =============================================================================
%  Spurious Precision: Why Shift-Share Designs Overstate Monetary Policy
%                      Effects on Local Employment
%  Plain-Language Summary + Introduction — draft of August 11, 2026
%  Author: Abdallah Dalis
%
%  Written after Results, Discussion and Methods, per step 6 of
%  00_Resources/research-paper-writing-process.md. Nothing here introduces a
%  claim those sections do not already support.
%
%  ⚠️ REFERENCES ARE NOT YET VERIFIED FOR THIS PAPER. Keys used below:
%     carlino1998         Carlino & DeFina. Corrected to 1998 (not 1999) during
%                         the Part 3 reference check, June 15 2026.
%     bartik1991, goldsmithpinkham2020, jorda2005, nakamura2018, bauer2023
%                         Carried over from the Part 3 bibliography, which was
%                         verified against publisher records June 15 2026.
%     adao2019            Adao, Kolesar & Morales, "Shift-Share Designs: Theory
%                         and Inference." PIPELINE_RECONCILIATION.md section 4b
%                         says VERIFY BEFORE CITING and that has not been done.
%                         Do not compile this to a submitted PDF until it is.
% =============================================================================

% -----------------------------------------------------------------------------
%  PLAIN-LANGUAGE SUMMARY BOX
%  Required on every empirical paper. Style: shaded box titled "Plain-Language
%  Summary", per Research_Monetary_Incidence/style_template/dalis-econ-style.sty.
%  Environment name below assumes that .sty; adjust if the macro differs.
% -----------------------------------------------------------------------------
\begin{plainsummary}

\textbf{What was asked.} When the Federal Reserve raises interest rates, some
local economies should suffer more than others. Places built on construction
and factories borrow more and build more. A standard method finds that these
places actually do \emph{better} when rates rise, and finds it with great
confidence. This paper asks whether that confidence is real.

\textbf{What was found.} It is not. The method measures policy using the
interest rate the Fed actually set. But the Fed raises rates when the economy is
already strong, so that measure carries the economy inside it. Swapping in a
measure of only the surprise part of policy leaves the estimates far too
uncertain to say anything. The margin of error grows by roughly five to twelve
times, depending on the sample.
Separately, the method fails its own built-in check for whether the comparison
was fair to begin with.

\textbf{What it means.} The tight margins in this kind of study can come from
economic ups and downs rather than from policy. Researchers using this design
should report both measures side by side and compare the margins of error, not
just the estimates. A large gap between them is a warning.

\textbf{The honest limitation.} This paper shows the original confidence was
misplaced. It does not recover the true effect. The cleaner estimates are too
imprecise to say whether employment rises or falls, and one further test of the
margins of error remains undone.

\end{plainsummary}


\section{Introduction}

% --- Establish the territory -------------------------------------------------
Monetary policy does not land evenly across a country. A rate rise raises the
cost of borrowing, and places whose employment sits in interest-sensitive
industries should feel it first. Construction and durable manufacturing are the
usual examples. \citet{carlino1998} established the regional version of this
claim, and it has been standard since.

The workhorse tool for testing it is the shift-share, or Bartik, design
\citep{bartik1991}. A national shock is interacted with a local exposure share
fixed before the sample begins. Variation in exposure does the identifying work,
and the fixed share is meant to keep exposure from responding to the shock
\citep{goldsmithpinkham2020}.

% --- Establish the niche -----------------------------------------------------
The design has a measurement problem at its centre, and it is not a subtle one.
Applied work usually measures policy with the realized change in the federal
funds rate. That change is endogenous. The Federal Reserve tightens when the
economy runs hot, and the sectors that define exposure are the most cyclical
ones in the economy.

Everyone writing in this literature knows that. The high-frequency
identification programme exists because of it \citep{nakamura2018, bauer2023}.
What is missing is a price. How much does the endogenous measure distort a
local labor market design, and along which dimension does the distortion show
up?

The instinctive answer is bias, meaning the coefficient sits in the wrong place.
This paper finds something different, and more awkward.

% --- Occupy the niche --------------------------------------------------------
I estimate a shift-share difference-in-differences design on county employment
from 2002 to 2019, using the predetermined 2002 county share of employment in
construction and manufacturing as exposure. Every result appears for the full
window and again excluding 2009 through 2012.

Three findings follow.

% RE-ESTIMATED August 12, 2026 on the rebuilt exposure measure. 0.033 supersedes
% 0.022. Source: all_coefficients_R_zlb-full_pos-off.csv, spec main_spec.
% The qualitative claim was unaffected, as predicted: the coefficient rose by
% about half, the standard error rose with it, and the t-statistic moved from
% 5.12 to 5.07.
First, the naive design looks strong and points the wrong way. The
contemporaneous interaction is $0.033$ with a $p$-value below $0.001$. Read
directly, tightening raises employment where exposure is highest.

% RESOLVED August 12, 2026. Count holds at two of four; the leads changed from
% {2,4} to {2,3}. This sentence does not name them, so only the count mattered.
Second, the design fails its own falsification test. Adding leads of the
interaction, two of the four reject zero under both clustered and randomization
inference. Employment in exposed counties was already moving before policy
changed.

% ⚠️ "every horizon turns negative" is a FULL-WINDOW claim only. See
% sections/results_discussion.tex 5.6. Do not restore it unqualified.
Third, and this is the paper's central result, replacing the realized rate
change with Bauer--Swanson high-frequency surprises does not relocate the
estimate. It dissolves it. No horizon rejects zero, in either window, and
standard errors widen by a factor of five to ten. The two estimators are not
statistically distinguishable at any horizon.

% --- The mechanism, briefly. Detail is in Results. ---------------------------
The mechanism is visible in the shock series themselves. Over this sample the
realized rate change correlates with contemporaneous national employment growth
at $0.470$. The identified surprise correlates at $-0.119$, and the two measures
correlate with each other at $-0.010$.

So the naive design draws its precision from business-cycle variation rather
than policy variation. Cyclical variation is abundant and moves with the
outcome, which is what produces tight intervals. It is also the variation an
identification strategy exists to remove. The design was not estimating the
effect of monetary policy precisely. It was estimating something else precisely.

% --- Contribution, stated narrowly on purpose --------------------------------
The contribution is diagnostic rather than structural. This paper does not
propose a new estimator, and it does not recover the true employment response.
It shows that a widely used design can report a precise and interpretable
coefficient whose precision is an artifact, and that two cheap checks catch it.

The recommendation follows from that and costs one column. Report the design
under both shock measures, and read the comparison on the standard errors
rather than the coefficients. A large gap says the reported precision was
borrowed.

% --- Honest placement relative to the shift-share inference literature -------
This is a complement to the shift-share inference literature, not a substitute
for it. \citet{adao2019} show that conventional standard errors can be severely
understated when a small number of common sectoral shocks drive exposure. That
concern is live here and is not tested in this paper. The problem documented
here is upstream of it: before asking whether the standard errors are computed
correctly, one may ask what variation they are computed on.

% --- Roadmap ----------------------------------------------------------------
Section 2 describes the data and both designs. Section 3 reports the results,
figure by figure. Section 4 discusses what the precision was made of, what the
two diagnostics do and do not establish, and what remains untested.
